% =========================
% report/experiments.tex
% =========================

\subsection{هدف آزمایش‌ها}
آزمایش‌ها برای پاسخ به سه سؤال طراحی می‌شوند:
\begin{enumerate}
  \item سامانه در شرایط معمول تا چه حد هم‌ترازی صحیح می‌دهد؟
  \item پایداری آن در برابر تغییرات نور، زاویه دید و شلوغی پس‌زمینه چقدر است؟
  \item هزینهٔ زمانی آن برای اجرا در کاربرد واقعی چه میزان است؟
\end{enumerate}

\subsection{پارامترهای کلیدی و بازه‌ها}
برای ارزیابی حساسیت و انتخاب رژیم پایدار، حداقل مجموعهٔ زیر پیشنهاد می‌شود:
\begin{itemize}
  \item تعداد ویژگی‌های ORB: \texttt{nfeatures} $\in \{500, 1000, 2000, 4000\}$،
  \item آستانهٔ RANSAC: $\tau \in \{2,3,5,8\}$ پیکسل،
  \item Cross-check: روشن/خاموش،
  \item CLAHE (اختیاری): روشن/خاموش.
\end{itemize}

\subsection{روش اجرای آزمایش}
برای هر زوج تصویر:
\begin{enumerate}
  \item اجرای پیش‌پردازش و استخراج ویژگی،
  \item تطبیق و پالایش تطبیق‌ها،
  \item تخمین $\mathbf{H}$ با RANSAC،
  \item تولید هم‌پوشانی و خروجی‌های دیباگ،
  \item ثبت معیارهای کمی در \texttt{per\_image.csv}.
\end{enumerate}

\subsection{خروجی‌های الزامی برای گزارش}
حداقل خروجی‌های زیر باید تولید شوند:
\begin{itemize}
  \item تصاویر هم‌پوشانی نمونه‌های نماینده،
  \item هیستوگرام خطای بازفرافکنی (اگر زمین‌حقیقت موجود است)،
  \item نمودار توزیع زمان اجرا (یا خلاصهٔ آماری مثل میانه و صدک ۹۵)،
  \item نمونه‌های شکست و تحلیل علت.
\end{itemize}
